\lecture{8}{Tue 31 May 2022 08:47}{Partially Adversarial SP with Examination}

\newpage
\section{Partially Adversarial SP}

\subsection{Problem formulation }
\textbf{Random arrival.} $n $ items are totally ordered. A random permutation $\sigma $ is taken uniformly from the permutation group $S_n$. The classical SP corresponds to finding a stopping time that most likely stops on $t = \sigma(1)$. We will always use numbers to denote the true rank of items, so the incoming stream of items is represented by $\sigma(1,\ldots,n)$.

\textbf{Adversarial noise.} The adversary adds a random noise with resultant permutation $\pi $. This permutation is not over the arriving order of object, but over the observed rank, i.e. $\pi(j)=i $ means that the item with actual rank $j $ will be found by the player to be $i $-th best, if the player observes all items without examining them.

\textbf{Restriction on adversary.} To restrict the power of adversary, we demand $\pi$ to be long-range stable, i.e. $\pi(i)<\pi(j)$ whenever $i<j $ and $j-i>c $, where the constant $c $ is known to the algorithm.

\textbf{Examination.} For each incoming item the player can choose from three options: \textsf{Accept}, \textsf{Decline}, and \textsf{Examine}. If the player chooses to \textsf{Examine}, the current item will be added to a set $E $ of examinated items, and the next item will be forfeited and left unobserved. We use $E_t$ to denote the items examined by the start of $t $-th round. The \textsf{Examine} option is not available if the current item is the last item.

\textbf{Observable quantities.} 
Let $O_t$ be the observed total order on round $t $. It is a total order on $[t] $ induced by $\pi \sigma^{-1}$.
Let $\widehat{O}_t$ be the actual total order on round $t$. It is a total order on $[t] $ induced by $\sigma^{-1}$.
Let $P_t $ be the total order obtained from examination, defined by restricting $\widehat{O}_t $ to $E_t $. Note that $O_t$ and $P_t$ may not be compatible since $O_t$ is subject to observational "mistakes" caused by the noise.
We use $\ge_{O_{t}}$ to denote the order inside the ordered set $O_{t}$.

It is equivalent, but sometimes more convenient, to use the notion of relative rank instead of total order. We define $\text{RR}_t(i)=\sum_{s=1}^t \mathbf{1}\left( \pi\sigma^{-1}(s)\ge \pi\sigma^{-1}(i) \right) $ to be the observed relative rank function.

\textbf{Decision rule.} Let $\left( \mathcal{F}_t \right)_{t \in [n]}$ be the sequence of $\sigma$-algebras with each $\mathcal{F}_t$ generated by the random variables $O_1 \times P_1, \ldots, O_t \times P_t$, that is, \[
\mathcal{F}_t = \sigma\left( O_1 \times P_1, \ldots, O_t \times P_t \right) =\sigma\left( O_t \times P_t \right) 
.\]  
By a stopping time, we mean a random variable $\tau$ taking values in $[n]$ and satisfying the property $\left\{ \tau=t  \right\} \in \mathcal{F}_t$. The decision rule is a stopping time $\tau$ together with a sequence of indicators $\textsf{Exm}_t :=\mathbf{1}(t \in E_{t+1})$ such that \[
	\textsf{Exm}_{t+1} \in \mathcal{F}_t \text{ and } \textsf{Exm}_{t}\textsf{Exm}_{t+1} = 0 \text{ and } \mathbf{1}\left( \tau = t  \right) \textsf{Exm}_t = 0
.\]

\textbf{Goal.} Our goal is to choose a decision rule $\left( \tau, \textsf{Exm} \right) $ to maximize $\operatorname{Pr}\left[{\tau=\sigma(1)}\right]$.

\subsection{Preliminaries}

	We list several obvious facts about this problem construction.
\begin{lemma}
\label{partial-adversarial-obv-facts}
	\begin{enumerate}
		\item The relative ranks behave the same:  $\text{RR}(t)$ is uniformly sampled from $1, \ldots,t+1$ and independent from each other.
		\item When $\pi$ do not depend on $\sigma$, $\sigma \pi$ and $\sigma$ have exactly same distribution.
	\end{enumerate}
\end{lemma}

\subsection{Case without examination}

We first do not allow examination, and see how far the algorithm can get.

\subsubsection{Weak adversarial, do not know anything}

Consider the weakest adversary, that can only pick a single $\pi$ without knowing about the algorithm. The competitive ratio is
\begin{align*}
	&\inf_\pi \sup_\tau \mathbb{E}_\sigma \left[ \operatorname{Pr}\left[{  \tau = \sigma\pi(1) }\right]   \right]\\
	&= \inf_k \inf_{\pi:\pi(1)=k}\sup_\tau \mathbb{E}_\sigma \left[ \operatorname{Pr}\left[{  \tau = \sigma(k) }\right]   \right]\\
	&\le  \inf_k \sup_\tau \mathbb{E}_\sigma \left[ \operatorname{Pr}\left[{  \tau = \sigma(k) }\right]   \right]
.\end{align*}

The problem reduces to a generalized post-doc variant of SP, where we want the item with absolute rank $k$ instead of $1$. Each time $k $ is provided, we may specifically design our algorithm to maximize the probability of choosing the rank $k $ item. This is significantly harder to solve as compared to the classical SP, in its value function is no longer monotone: the function \[
\operatorname{Pr}\left[{\text{AR}(t) = m | \text{RR}(t) = r }\right] = \frac{\binom{m-1}{r-1}\binom{n-m}{t-r}}{\binom{n}{t}}
\]
is not monotone wrt $r $ for any fixed $m $. No explicit solution has been found so far. We can show that the worst case when $c < \frac{n}{2}$ is attained at $\pi(1)=c $, and at $\pi(1) = \left\lfloor \frac{n}{2} \right\rfloor $ when $ c \ge  \frac{n}{2}$.

\subsubsection{Weak adversarial, know algorithm}

Now we assume that the adversary adapts to each algorithm, but cannot depend on the random permutation $\sigma$. The competitive ratio is now
\begin{align*}
	&\sup_\tau  \inf_\pi \mathbb{E}_\sigma \left[ \operatorname{Pr}\left[{ 
	\tau = \sigma\pi(1) }   \right] \right]\\
\intertext{Note that the generalized postdoc variant of SP is an upper bound to the problem:}
	&= \sup_\tau \inf_k \inf_{\pi:\pi(1)=k } \mathbb{E}_\sigma \operatorname{Pr}\left[{ \tau =\sigma(k) }\right]  \\
	&\le  \sup_\tau \inf_k \mathbb{E}_\sigma \operatorname{Pr}\left[{ \tau =\sigma(k) }\right]  
.\end{align*}
This time we need to tackle all sub-problems using only one algorithm. In other words, the algorithm's probability to choose the 1, 2, \ldots, $c $-th item are compared, and the minimum one is the CR.

For lower bound, we consider the following algorithm: When each item arrives, if $\text{RR}(t) > c $, reject the item. Otherwise, accept with probability $p $, where $p $ is a tuning parameter. Now for some algorithm $\tau$ and fixed $k $, We derive analytic result for this algorithm. Its performance monotonously decrease wrt $\pi(1)$ for fixed $n $ and $p $, and always approach the trivial lower bound $\frac{1}{n }$. 

Consider another algorithm: Let the first $pn$ items to go by, and then select the first item with relative rank better or equal to $c $. This algorithm makes no distinction among the first $c $ items, so it suffices to analyze its behavior on the classic SP. 

\subsubsection{Strong adversarial, knows $\tau $ and $\sigma$}

\begin{align*}
	\sup_\tau  \mathbb{E}_\sigma \inf_\pi \left[ \operatorname{Pr}\left[{ 
	\tau = \sigma\pi(1) }   \right]\right] 
.\end{align*}





\subsubsection{Strong adversarial, knows $\tau$,  $\sigma$ and $\mathbf{r}$ }


\begin{align*}
	\sup_\tau  \mathbb{E}_\sigma \left[ \mathbb{E}_{\mathbf{r}} \left[ \inf_\pi  
	\mathbf{1}\left( \tau = \sigma\pi(1) \right)  \right]  \right] 
.\end{align*}

In this case the adversary knows the random seed used by the algorithm, so it's able to make hindsight decisions. In this setting, the algorithm achieve zero cr. This can be shown by calling to an auxiliary problem of adversarial best-choice within the top $c $ items, where no deterministic algorithm can guarantee nonzero competitive ratio against an adversary.


\subsection{Lower Bound for an algorithm}

We establish a lower bound for the following simple algorithm: examine each item with $\text{RR}_t(t) \le  c $. Accept the current item if $t \ge  T $ and the item is the best so far.

\begin{algorithm}[H]
\caption{Simple algorithm}
\label{alg:simple} 
\begin{algorithmic}[1]
\While{At least two items are left}
\State{$x \gets$ the current item to be decided}
\State{$t, O_{t}$, $P_{t} \gets $ as defined before }
\If{$x \in P_{t}$}
\If{$x = \max P_{t}$ and $t\ge T$}
\State{\textsf{Accept} $x$}
\Else
\State{\textsf{Decline} $x$}
\EndIf
\ElsIf{$\left| \{a \in O_{t}: a \ge_{O_t} x\} \right| \le c $} \Comment{Means $\text{RR}_t(x)\le c$}
\State{\textsf{Examine} $x$}
\Else
\State{\textsf{Decline} $x$}
\EndIf
\EndWhile
\State{\textsf{Accept} the only item left}
\end{algorithmic}
\end{algorithm}

\begin{lemma}[Equivalent construction of the game]
It is equivalent to view this game as follows: when an item is examined, instead of omitting the next item, the last available item is taken from the end of the queue. The player wins only if the best item is accepted, and the item is not omitted.
\end{lemma}
\begin{proof}
	(informal, to be put rigorously) Consider a deterministic algorithm. The algorithm never knows $\sigma^{-1}(t)$, the position of the current item in the arrival sequence. For some $\sigma \in S_n$, we can perform some reordering according to the following rule: We first let the algorithm run through the sequence, and take out all the items omitted, and stack them to the end of the queue. The resultant sequence $\sigma'$ looks exactly same to the algorithm in the new setting.
	Hence, the whole gameplay will be identical up to an automorphism on $S_n$, and there's no difference revealed to the algorithm. Because the only randomness we used is the randomized $\sigma \in  S_n$, we have essentially created an isomorphism between the two probability spaces for the two versions of the game.
\end{proof}

\begin{intuition}
	Imagine a strip of paper which represents the incoming sequence of items. One consumes the items from left to right, but the paper strip burns from right to left (due to examination which makes some items unavailable). The moment the paper is used up, the game is forced to end.
\end{intuition}
\begin{figure}[ht]
    \centering
    \incfig{lowerbound}
    \caption{Our scheme to finding a lower bound}
    \label{fig:lowerbound}
\end{figure}

Let $t^* := \sigma^{-1}(1)$, i.e., the position of the best item. We couple the game with the classical secretary problem (CSP). 

\begin{lemma}[Coupling with CSP]
\label{lem:CSP-couple}
	Whenever an algorithm wins the problem, it automatically wins CSP in the same item setting. Whenever an algorthm wins CSP such that $t^* + |E_{t^*}| \le n$, it automatically wins our problem.
\end{lemma}
\begin{proof}
	To be done.
\end{proof}

\begin{lemma}
	\label{lem:CSP*}
	For any $a$ and $n> a$,
 	\begin{equation*}
		\begin{split}
	&\operatorname{Pr}\left[{\text{win CSP with }a \text{ items } |  t^* +  |E_{t^*}| \le n }\right] \\
	&\qquad \ge \operatorname{Pr}\left[{\text{win CSP with }a \text{ items }}\right]
		\end{split}
 	\end{equation*}  

\end{lemma}
\begin{proof}
	To be done.
\end{proof}


From Lemma \ref{lem:CSP-couple} we have
\begin{align}
	 \operatorname{Pr}\left[{\text{win}}\right] 
	 &= \operatorname{Pr}\left[{\text{win CSP with }n \text{ items and } t^* +  |E_{t^*}| \le n}\right]\\
	 \intertext{Now we constrain $t^*$ and $|E_{t^*}|$ by  $a$ to make sure they don't meet.}
	   &\ge \max_{a \le n} \operatorname{Pr}\left[{\text{win CSP with }n \text{ items and } t^* +  |E_{t^*}| \le n \cap t^* \le a}\right]\\
	 \intertext{Which follows from property of intersection. Continue to apply elementary properties,}
	 \begin{split}
	   &= \max_{a \le n} \operatorname{Pr}\left[{\text{win CSP with }n \text{ items } |  t^* +  |E_{t^*}| \le n \cap t^* \le a}\right] \\
	   & \qquad\qquad\qquad \operatorname{Pr}\left[{t^* +  |E_{t^*}| \le n | t^* \le a}\right]\operatorname{Pr}\left[{t^* \le a}\right]
	 \end{split}\\
	 \begin{split}
	   &= \max_{a \le n} \operatorname{Pr}\left[{\text{win CSP with }n \text{ items } |  t^* +  |E_{t^*}| \le n \cap t^* \le a}\right] \\
	   &\qquad\qquad\qquad \operatorname{Pr}\left[{|E_{t^*}|\le n-a}\right] \operatorname{Pr}\left[{t^* \le a}\right]
	 \end{split}\\
	 \intertext{Since a CSP with n items where only the first $a$ items may be the best is identical to a CSP with $a$ items,}
	 \begin{split}
	   &= \max_{a \le n} \operatorname{Pr}\left[{\text{win CSP with }a \text{ items } |  t^* +  |E_{t^*}| \le n}\right] \\
	   &\qquad\qquad\qquad \operatorname{Pr}\left[{|E_{t^*}|\le n-a}\right] \operatorname{Pr}\left[{t^* \le a}\right]
	 \end{split}\\
	 \intertext{By Lemma \ref{lem:CSP*}}
	 \begin{split}
	   &= \max_{a \le n} \operatorname{Pr}\left[{\text{win CSP with }a \text{ items } }\right] \\
	   &\qquad\qquad\qquad \operatorname{Pr}\left[{|E_{t^*}|\le n-a}\right] \operatorname{Pr}\left[{t^* \le a}\right]\\
	 \end{split}\\
	 \intertext{Evaluate each term separately,}
	 \operatorname{Pr}\left[{\text{win}}\right] 
	   &\ge \max_{a \le n} \frac{1}{e } \operatorname{Pr}\left[{|E_a|\le n-a}\right] \operatorname{Pr}\left[{t^*\le a}\right]
.\end{align}

For $t\ge c$ we have \[
|E_t| = c-1+c \sum_{s=c}^{t} B_s
\]  where $B_s \sim \text{Ber}(\frac{1}{s })$.

Using Hoeffding's Bound:

\[
\operatorname{Pr}\left[{\sum_c^t B_s - \sum_c^t \frac{1}{s } \le \beta}\right] \ge 1-\exp\left( - \frac{2\beta^2}{t-c+1} \right) 
\] 
\[
\operatorname{Pr}\left[{|E_t| \le c(\beta+\sum_{s=c }^t \frac{1}{s })+c-1}\right]\ge 1-\exp\left( - \frac{2\beta^2}{t-c+1} \right) 
.\] 

Fix some $a$. Let $\beta$ be the greatest real number such that \[
c(\beta+\sum_{s=c }^a \frac{1}{s })+c-1 \le n-a
.\] 
Therefore,  \[
\beta = \frac{n-a}{c} -1 - H(a)+H(c-1)
\] where we use $H(a)$ to denote the $a $-th harmonic number.

Now we have a concentration bound for $|E_{t^*}|$ :
\[
	\operatorname{Pr}\left[{|E_{t^*}| \le n-a}\right]\ge 1-\exp\left( - \frac{2\beta^2}{t-c+1} \right) 
.\] 

Applying this bound, we now have
\begin{align*}
	\operatorname{Pr}\left[{\text{win}}\right] \ge  \max_{a\le n } \frac{1}{e } \frac{a}{n } \left( 1-\exp\left( - \frac{2\beta^2}{a-c+1} \right)  \right) 
.\end{align*}


\textbf{Numerical Result.}
Take $c=2 $ and $n = 100$ for an example. When $a = 70$, we have the numerical bound $\operatorname{Pr}\left[{\text{win }}\right]\ge 0.674 \frac{1}{e}$, which is a non-trivial improvement from the examine-everything algorithm which guarantees only a competitive ratio of $0.5 \frac{1}{e}$.

Take $c = 1$ and $n = 100$, then the optimal choice of $a$ is $82$, and the corresponding bound is $\operatorname{Pr} \left[\text{win}\right] \ge 0.807 \frac{1}{e}$.

\newpage
\subsection{Upperbound for the problem}

\textbf{This section only contains ideas, there's not any proof.}

\subsubsection{some guesses}

\begin{definition}
	A \textit{simple algorithm} is an algorithm in the following form, with parameters $c$ and $T$ replaced by any valid value: \textsf{Examine} each item with observed relative rank $\le c$, and \textsf{Accept} if it is $\max P_{t+1}$ and  $t\ge T$.
\end{definition}

\begin{conjecture}
	Any algorithm, no matter deterministic or randomized, has its competitive ratio dominated by some simple algorithm.
\end{conjecture}

\begin{conjecture}
	For each $c$ and $n$ (as parameters of the game), of all simple best algorithms, there is one optimal simple algorithm. The optimal algorithm has the exactly same $c$.
\end{conjecture}

If the above guesses are true, then we only need to find the optimal $T$ for the simple algorithm, and bound its competitive ratio from above.


\subsubsection{An alternative thought}

The process of examination and omitting items will often bring the player to a \textbf{from-two-choose-one} scenario. For example, consider the following sequence of items:

$*, \triangleleft , \triangleright , -, -$

Where the player is at *. Suppose that both $\triangleleft$  and $\triangleright$  are "critical" items such that they have relative rank $\le $ c and shall be examinated. Now if the player declines *, it has to examine $\triangleleft$  and therefore lose $\triangleright$ ; if it examines and declines *, it loses $\triangleleft$ .

Therefore one of $\triangleleft$  and $\triangleright$  must be lost, and there is no way to mitigate this.
This behavior can be something fundamental that can be properly exploited to obtain an upper bound.
